% ============================================================
% A Spherical N-Series Area Law for Regular Refinement Families
% Version 2.3 release candidate -- October 29, 2026
%
% Editorial provenance notes (the unresolved items are retained as limitations):
%  * Every item that could not be reconstructed from the v2.1 text is
%    described in visible provenance notes, not represented as verified runs.
%  * Constants -1/6, -1/90, -1/810 were verified two ways:
%    (a) direct summation of the weights in Eqs. (2)-(4);
%    (b) product of Romberg stage factors (4^k 2^{-3}-1)/(4^k-1).
%  * Numerical content of Tables 1-8 is copied unchanged from v2.1.
% ============================================================
\documentclass[11pt]{article}
\usepackage[utf8]{inputenc}
\usepackage[T1]{fontenc}
\usepackage{amsmath,amssymb,amsthm}
\usepackage{booktabs}
\usepackage{xcolor}
\usepackage[margin=1in]{geometry}
\usepackage{hyperref}


\newtheorem{proposition}{Proposition}
\theoremstyle{remark}
\newtheorem{remark}{Remark}

\newcommand{\Ahat}{\widehat{A}}
\newcommand{\Aex}{A_{\mathrm{exact}}}

\title{A Spherical N-Series Area Law for Regular Refinement Families\\[2pt]
\large Analytic regular-family expansions and conditional perturbation law}
\author{C. Wayne Baker}
\date{Version 2.3 curated release candidate --- October 29, 2026\\ Original manuscript: June 3, 2026}

\begin{document}
\maketitle

\begin{abstract}
We study multiscale extrapolation of spherical areas computed from geodesic boundary
polygons and chordal surface triangulations. The previously reported numerical experiments
show observed orders close to $2,4,6,8$ for spherical caps, normalized-barycentric spherical
triangles, latitude--longitude cells, and regular chordal refinements of fixed irregular
polygons. For spherical caps and uniform latitude--longitude chordal cells away from the poles,
we justify the even-power asymptotic structure analytically; other reported geometries remain
numerical examples. The underlying scale weights are the classical Richardson--Romberg weights
specialized to even error exponents. This revision sharpens the regularity hypothesis and
analyzes a possible reason that the high-order ladder deteriorates under scale-dependent
interior-node perturbation. Under an explicit smooth-coefficient expansion, displacements
proportional to the grid spacing can generate an odd $h^3$ error term that survives the
even-power extrapolation; we give its exact residual factors ($-1/6$, $-1/90$, $-1/810$ for
the two-, three- and four-scale operators). We give exact generalized scale weights that can
instead eliminate a prescribed set of even and odd exponents. The new perturbation mechanism
is a conditional mathematical explanation, and the exploratory AI-assisted fit diagnostics, which lack complete provenance, are excluded from this version. No general
even-power theorem for arbitrary spherical meshes is claimed.
\end{abstract}

%------------------------------------------------------------
\section{Introduction}
Many geometric area approximations admit structured error expansions in a refinement
parameter $h\asymp N^{-1}$. When the errors are in even powers, the same approximation
evaluated at successive scales can be combined to remove those terms. In the spherical
cases studied here, the observed sequence is
\[
2\longrightarrow 4\longrightarrow 6\longrightarrow 8 .
\]
We examine (i) a spherical cap bounded by a regular geodesic polygon, (ii) a chordally
triangulated birectangular spherical triangle, (iii) a chordally triangulated
latitude--longitude cell, (iv) fixed irregular spherical polygons under regular chordal-fan
refinement, and (v) perturbed interior grids. The original tabulated experiments are retained below with explicit provenance limitations.
New Sections~4 and~9 distinguish geometric nonuniformity from an approximation family whose
definition changes with the refinement scale.

%------------------------------------------------------------
\section{Contributions and scope}
The principal empirical observation is that several geometrically distinct, fixed refinement
constructions behave as if
\begin{equation}\label{eq:expansion}
A_N(D)=A(D)+C_2(D)N^{-2}+C_4(D)N^{-4}+C_6(D)N^{-6}+\cdots .
\end{equation}
For families actually possessing this expansion through the required order, extrapolation
raises the convergence order successively from 2 to 4, 6, and 8. A direct analytic argument
establishes the even-power structure for the cap construction. The other original experiments
provide numerical evidence, not a universal proof. Fixed irregular boundaries need not destroy
the ladder; in contrast, changing a mesh-generating map with $N$ may introduce uncancelled odd
powers. The operator mechanism itself is classical; the geometric question is which families
satisfy its asymptotic hypotheses.

%------------------------------------------------------------
\section{N-Series scale operators}
Put $h=N^{-1}$, let $TA(h)=A(h/2)$, and set $A_N=A(h)$. If
$A(h)=A+c_2h^2+c_4h^4+c_6h^6+O(h^8)$, define
\begin{align}
\Ahat^{(2)}_N &= \frac{4A_{2N}-A_N}{3}, \label{eq:R2}\\
\Ahat^{(3)}_N &= \frac{A_N}{45}-\frac{4A_{2N}}{9}+\frac{64A_{4N}}{45}, \label{eq:R3}\\
\Ahat^{(4)}_N &= -\frac{A_N}{2835}+\frac{4A_{2N}}{135}-\frac{64A_{4N}}{135}+\frac{4096A_{8N}}{2835}. \label{eq:R4}
\end{align}
Each set of weights sums to one and eliminates $h^2$ and, respectively, the additional powers
$h^4$ and $h^6$. Hence, subject to a sufficiently long asymptotic expansion and nonzero next
coefficient, the errors are of orders 4, 6, and 8. These are Richardson--Romberg extrapolation
formulas in a geometric N-Series setting.

In terms of the shift $T$ (which commutes with itself and acts on a pure power $h^p$ as
multiplication by $2^{-p}$), Eqs.~\eqref{eq:R2}--\eqref{eq:R4} are the nested Romberg stages
\[
\Ahat^{(r)}=\prod_{k=1}^{r-1}\frac{4^{k}T-I}{4^{k}-1}\,A,\qquad r=2,3,4 .
\]
For example, $(16T-I)(4T-I)/45=(64T^2-20T+I)/45$ reproduces Eq.~\eqref{eq:R3}. On an $h^p$
term, stage $k$ contributes the factor $(4^k2^{-p}-1)/(4^k-1)$, which vanishes exactly when
$p=2k$. This product form gives the residual moments used in Section~9.

%------------------------------------------------------------
\section{An accurate regularity statement}
For a fixed spherical domain $D$ and a specified geometrical area rule, call a refinement
family \emph{even-asymptotic through order $2r$} when, as $h\downarrow0$ along the admissible
refinement sequence,
\begin{equation}\label{eq:evenasym}
A(h)-A(D)=\sum_{k=1}^{r}c_{2k}h^{2k}+O(h^{2r+2}),
\end{equation}
with coefficients independent of the refinement level. Equation~\eqref{eq:evenasym}, rather
than informal smoothness alone, is the operative hypothesis of the $2,4,6,\dots$ ladder. For a
given scheme, a fixed sufficiently smooth reparametrization of a reference mesh is a natural
candidate for a scale-consistent nonuniform family: the map is independent of $h$, while only
the uniform reference-grid spacing changes. However, being fixed and smooth is not on its own a
theorem that every area discretization satisfies \eqref{eq:evenasym}. Even-power structure must
be proved for the particular geometry and area rule or supported by convergence tests. The map
should also preserve orientation and describe the same target domain at every scale.

%------------------------------------------------------------
\section{Evidence I: spherical cap area}
For a cap with polar angle $\alpha$ on the unit sphere, $A_{\mathrm{cap}}=2\pi(1-\cos\alpha)$.
Join the north pole to adjacent vertices of the regular geodesic boundary polygon. With
$c=\cos\alpha$, $s=\sin\alpha$, and $h=2\pi/N$, the oriented spherical triangle area for
sufficiently small $h>0$ is
\begin{equation}\label{eq:E}
E_\alpha(h)=2\arctan\!\left(\frac{s^{2}\sin h}{1+2c+c^{2}+s^{2}\cos h}\right).
\end{equation}
(This follows from $\tan(E/2)=t\sin h/(1+t\cos h)$ with $t=\tan^2(\alpha/2)=(1-c)/(1+c)$,
after multiplying numerator and denominator by $(1+c)^2/(1+c)$ and using $1-c^2=s^2$.)
For $0<\alpha<\pi$ the denominator is positive near zero and even in $h$, while the numerator
is odd. Therefore $E_\alpha$ is odd analytic near zero:
\[
E_\alpha(h)=a_1h+a_3h^3+a_5h^5+\cdots,\qquad a_1=1-\cos\alpha .
\]
Since $A_N=NE_\alpha(2\pi/N)$, the polygon area has an even-power error expansion, and its
constant term $2\pi a_1=2\pi(1-\cos\alpha)$ equals the cap area exactly. This is a proof for
this cap family, not for every triangulated sphere.

\begin{table}[ht]
\centering
\begin{tabular}{lrrrr}
\toprule
$\alpha$ & Raw & Two-scale & Three-scale & Four-scale\\
\midrule
$15^\circ$ & 1.99989583 & 3.99999065 & 6.00093925 & 7.99993479\\
$30^\circ$ & 1.99996740 & 4.00041698 & 6.00000053 & 7.95025527\\
$60^\circ$ & 2.00016296 & 4.00014936 & 6.00003853 & 7.99880645\\
$75^\circ$ & 2.00023455 & 4.00028432 & 6.00028420 & 8.00026989\\
\bottomrule
\end{tabular}
\caption{Original cap experiment: observed orders formed by comparing $N=128$ with $N=256$ (the row label is the cap polar angle $\alpha$). The $30^\circ$ four-scale value $7.95025527$ is reproduced at 80-digit precision and represents a pre-asymptotic deviation, not a precision floor.}
\end{table}

%------------------------------------------------------------
\section{Evidence II: chordal spherical triangle area}
Take $A=(0,0,1)$, $B=(1,0,0)$, and $C=(\cos\beta,\sin\beta,0)$. This is a birectangular
spherical triangle: the sides $AB$ and $AC$ are meridians and $BC$ lies on the equator, so the
angles at $B$ and $C$ are right angles and the angle at the pole $A$ is $\beta$. Its spherical
excess is $\beta+\pi/2+\pi/2-\pi=\beta$, so $\Aex=\beta$. Map a normalized barycentric
triangulation to the sphere and sum the Euclidean areas of its chordal triangles.

\begin{table}[ht]
\centering
\begin{tabular}{lrrrr}
\toprule
$\beta$ & Raw & Two-scale & Three-scale & Four-scale\\
\midrule
$15^\circ$ & 1.99942072 & 3.99915618 & 5.99886117 & 7.99716985\\
$30^\circ$ & 1.99938698 & 3.99908575 & 5.99880136 & 7.99761444\\
$60^\circ$ & 1.99923969 & 3.99881849 & 5.99858915 & 7.99835976\\
$90^\circ$ & 1.99885313 & 3.99827889 & 5.99785193 & 7.99808456\\
\bottomrule
\end{tabular}
\caption{Original chordal triangle experiment: observed orders at $N=64$.}
\end{table}

%------------------------------------------------------------
\section{Evidence III: latitude--longitude cell area}
For the unit-sphere cell $\lambda_0\le\lambda\le\lambda_1$ and $\phi_0\le\phi\le\phi_1$, the
exact area is
\[
A_{\mathrm{cell}}=(\lambda_1-\lambda_0)(\sin\phi_1-\sin\phi_0).
\]
The original approximation uses a uniform latitude--longitude parameter grid and two Euclidean
chordal triangles per parameter cell.

\begin{table}[ht]
\centering
\begin{tabular}{lrrrr}
\toprule
Cell & Raw & Two-scale & Three-scale & Four-scale\\
\midrule
Equatorial $30^\circ\times30^\circ$ & 1.99996836 & 3.99999637 & 5.99994223 & 7.99911861\\
Midlatitude $30^\circ\times30^\circ$ & 1.99998486 & 3.99996146 & 5.99998582 & 8.00085300\\
Wide equatorial $60^\circ\times30^\circ$ & 1.99992674 & 3.99992210 & 6.00000864 & 8.00036842\\
Northern $15^\circ\times60^\circ$ & 1.99994200 & 3.99996772 & 5.99996443 & 8.00007636\\
\bottomrule
\end{tabular}
\caption{Original latitude--longitude experiment: observed orders at $N=64$.}
\end{table}


\begin{proposition}[Even-power expansion for uniform latitude--longitude chordal cells]
\label{prop:latlon-even}
Fix $-\pi/2<\phi_0<\phi_1<\pi/2$ and $\lambda_0<\lambda_1$.
Subdivide the corresponding unit-sphere parameter rectangle into a uniform
$N\times N$ grid, with every cell triangulated along its $(p_{00},p_{11})$
diagonal, and let $A_N$ be the sum of the unsigned Euclidean areas of its two
chordal triangles. With $h=1/N$, for every positive integer $r$ there are
coefficients $c_{2},\ldots,c_{2r}$, independent of $N$, such that
\begin{equation}\label{eq:latlon-even}
 A_N=(\lambda_1-\lambda_0)(\sin\phi_1-\sin\phi_0)
       +\sum_{k=1}^{r}c_{2k}h^{2k}+O(h^{2r+2}).
\end{equation}
The assertion concerns this fixed rectangular domain and uniform, consistently
triangulated refinement; it does not extend to arbitrary perturbed grids.
\end{proposition}
\begin{proof}
Put $\Delta\phi=\phi_1-\phi_0$ and $\Delta\lambda=\lambda_1-\lambda_0$.
A parameter cell has centre $(x,y)$ and corners
$(x\pm\Delta\phi\,h/2,y\pm\Delta\lambda\,h/2)$.
For the fixed $p_{00}$--$p_{11}$ diagonal, let $g(x,y,h)$ be the sum of the
unsigned Euclidean areas of the chordal triangles
$(p_{00},p_{10},p_{11})$ and $(p_{00},p_{11},p_{01})$.
Changing the sign of $h$ permutes the opposite corners; taken together, the
two unsigned triangle areas are invariant under this relabelling. Their sum therefore satisfies
$g(x,y,-h)=g(x,y,h)$ for real $h$ sufficiently close to zero.
Each triangle has area one half of the norm of a cross product of edge
vectors analytic in $h$. After factoring $h^2$, each squared norm has a
positive leading coefficient proportional to
$(\Delta\phi\,\Delta\lambda\cos x)^2$; hence its positive square root
is real analytic for small signed $h$, uniformly on the closed latitude band. Thus $g/h^2$ extends to an even real-analytic
function of $h$, uniformly with respect to $(x,y)$ on the fixed rectangle.
In particular, for any fixed $r$,
\[
 g(x,y,h)=\sum_{k=1}^{r+1}\rho_{2k}(x,y)h^{2k}
                +O(h^{2r+4}),\qquad
 \rho_2(x,y)=\Delta\phi\,\Delta\lambda\cos x,
\]
with a uniform remainder. Rotational invariance makes the coefficient
functions independent of longitude $y$, although this is not required for
the midpoint argument.
Summing the remainder over $N^2=h^{-2}$ cells gives
$O(h^{2r+2})$. For each coefficient, the cell-centre sum is a composite
midpoint quadrature on a fixed rectangle, and the Euler--Maclaurin formula
for smooth coefficient functions has an expansion in even powers of $h$.
Collecting terms yields \eqref{eq:latlon-even}. The leading midpoint integral
is
\[
 \int_{\phi_0}^{\phi_1}\int_{\lambda_0}^{\lambda_1}
            \cos\phi\,d\lambda\,d\phi
 = (\lambda_1-\lambda_0)(\sin\phi_1-\sin\phi_0).
\]
\end{proof}

\begin{remark}[Scope of the parity argument]
The midpoint nodes and the fixed diagonal are part of the definition of the
refinement family. The assertion is an asymptotic expansion to each finite
order, not a claim that an infinite series converges. In particular it does
not cover the scale-dependent interior-node perturbation in Section~9.
\end{remark}

%------------------------------------------------------------
\section{Evidence IV: irregular spherical polygon domains}
The original study considered fixed irregular convex spherical polygons triangulated from a
root vertex. Each fan triangle was refined using a normalized barycentric chordal patch. One
six-vertex example had angular vertices $(4^\circ,-8^\circ)$, $(18^\circ,5^\circ)$,
$(15^\circ,28^\circ)$, $(2^\circ,42^\circ)$, $(-14^\circ,30^\circ)$, and $(-10^\circ,4^\circ)$;
the reported reference area by spherical excess was
$0.325062777898240815426301649532$. The main inference is empirical: an irregular but fixed
domain boundary does not automatically prevent a clean ladder. It is important that the fan
geometry remains well posed, consistently oriented, and convergent to the intended spherical
region.

\begin{table}[ht]
\centering
\begin{tabular}{rrrrr}
\toprule
$N$ & Raw & Two-scale & Three-scale & Four-scale\\
\midrule
16 & 1.99929248 & 3.99269841 & 5.99184034 & 8.01301620\\
32 & 1.99982228 & 3.99817005 & 5.99798552 & 8.00310155\\
64 & 1.99995552 & 3.99954224 & 5.99949790 & 8.00076594\\
\bottomrule
\end{tabular}
\caption{Original fixed six-vertex polygon experiment.}
\end{table}

\begin{table}[ht]
\centering
\begin{tabular}{rrrrr}
\toprule
Vertices & Raw & Two-scale & Three-scale & Four-scale\\
\midrule
5 & 1.99993053 & 3.99981653 & 5.99991832 & 7.99955672\\
6 & 1.99993586 & 3.99979158 & 5.99991390 & 8.00029966\\
8 & 2.00001032 & 4.00140962 & 5.99925770 & 7.99912295\\
\bottomrule
\end{tabular}
\caption{Original aggregate final-level orders over deterministic polygon families and seeds.}
\end{table}

%------------------------------------------------------------
\section{Evidence V: perturbations and a testable mechanism}
In the original perturbation experiment, boundary nodes of a latitude--longitude cell remained
fixed, while interior nodes were displaced by a deterministic smooth field. Perturbation
amplitude was specified as a fraction of local grid spacing; three seeds were used at each
level.
Because the amplitude is a fraction of the local spacing, this experiment belongs to the
scale-dependent family discussed in Section~9.1, not to the fixed-map family of Section~4.

\begin{table}[ht]
\centering
\begin{tabular}{lrrrr}
\toprule
Relative perturbation & Raw & Two-scale & Three-scale & Four-scale\\
\midrule
0.00 & 1.99996836 & 3.99999637 & 5.99994223 & 7.99911861\\
0.02 & 1.99997529 & 4.02004815 & 4.53899093 & 2.61899119\\
0.05 & 2.00001207 & 3.44420303 & 5.68235637 & 2.77050496\\
0.10 & 2.00014360 & 3.89815337 & 8.94059560 & 3.59334301\\
0.20 & 2.00066838 & 3.90830154 & 6.15480240 & 4.36890149\\
\bottomrule
\end{tabular}
\caption{Original perturbed-grid measurements (Table~6 of the June 2026 manuscript):
aggregate final-level mean observed orders. These are diagnostics, not verified asymptotic
exponents.}
\end{table}

An apparent order such as $8.94059560$ is not evidence of asymptotic convergence above the
theoretical order: accidental coefficient cancellation, small denominators in measured error
ratios, sign changes, or numerical floors can distort finite-level slopes. The original data
establish only that the second-order raw behavior is comparatively robust and that
higher-order extrapolation is not uniform in this test window.

\paragraph{Qualitative reading of Table~6.}
At small perturbation the two-scale order stays near 4 (0.02: 4.02) while the three- and
four-scale orders fall well below their nominal 6 and 8, with four-scale values between about
2.6 and 4.4. This is qualitatively what a residual $h^3$ term would produce: the higher-order
operators remove the $h^4,h^6$ terms that otherwise mask it, and the residual factors of
Proposition~\ref{prop:cubic} below (e.g.\ $-1/6$ versus $-1/810$) are applied to a coefficient that is linear
in the amplitude. This is compatibility, not confirmation. The non-monotone entries, the
erratic three-scale values, and the demonstrable pre-asymptotic variation in the unperturbed
four-scale baseline (cf.\ Table~1, $30^\circ$) mean that the table cannot by itself separate a
cubic residual from finite-window or round-off effects.


%------------------------------------------------------------
\subsection{Conditional appearance of a cubic term}
The key distinction is between a fixed nonuniform parametrization and a parametrization whose
displacement amplitude changes with $h$. Suppose $A(h,\delta)$ is an approximation defined
using a perturbation parameter $\delta$, and the following smooth, uniform in $\delta$
even-power expansion holds near $(h,\delta)=(0,0)$:
\begin{equation}\label{eq:joint}
A(h,\delta)-A=C_2(\delta)h^2+C_4(\delta)h^4+O(h^6),
\end{equation}
where $C_2$ is three times continuously differentiable (or has a locally Lipschitz second
derivative) and $C_4$ is at least once continuously differentiable near zero. This is a
hypothesis for a particular numerical construction, not something automatically implied by
smooth mesh perturbations.

\begin{proposition}[Scale-dependent perturbations can create odd powers]\label{prop:cubic}
Assume \eqref{eq:joint} with a remainder uniform for $|\delta|\le\delta_0$ and set
$\delta=\varepsilon h$, with fixed $\varepsilon$. Then
\begin{equation}\label{eq:cubic}
A(h,\varepsilon h)-A=C_2(0)h^2+\varepsilon C_2'(0)h^3
+\Big[\tfrac12\varepsilon^2C_2''(0)+C_4(0)\Big]h^4+O(h^5).
\end{equation}
If $\varepsilon C_2'(0)\neq0$, then the two-, three- and four-scale even-power operators
\eqref{eq:R2}--\eqref{eq:R4} each leave a cubic residual
\[
-\tfrac16\,\varepsilon C_2'(0)\,h^3,\qquad
-\tfrac1{90}\,\varepsilon C_2'(0)\,h^3,\qquad
-\tfrac1{810}\,\varepsilon C_2'(0)\,h^3,
\]
respectively, so the extrapolated error is $O(h^3)$ and is not improved beyond third order by
adding scales.
\end{proposition}

\begin{proof}
Taylor-expand $C_2(\varepsilon h)$ and $C_4(\varepsilon h)$ in \eqref{eq:joint}; the $h^3$
coefficient is $\varepsilon C_2'(0)$. For a scale-doubling weight set $\{w_j\}$ the surviving
coefficient is $\varepsilon C_2'(0)\sum_jw_j2^{-3j}$. Direct summation gives
\begin{align*}
\text{(2)}&:\ -\tfrac13+\tfrac43\cdot\tfrac18=-\tfrac16,\\
\text{(3)}&:\ \tfrac1{45}-\tfrac49\cdot\tfrac18+\tfrac{64}{45}\cdot\tfrac1{64}=-\tfrac1{90},\\
\text{(4)}&:\ -\tfrac1{2835}+\tfrac4{135}\cdot\tfrac18-\tfrac{64}{135}\cdot\tfrac1{64}
+\tfrac{4096}{2835}\cdot\tfrac1{512}=-\tfrac1{810}.
\end{align*}
Equivalently, by the product form in Section~3 the stage factors at $p=3$ are $-\tfrac16$, $\tfrac1{15}$,
$\tfrac19$, whose cumulative products are $-\tfrac16,-\tfrac1{90},-\tfrac1{810}$.
\end{proof}

\begin{remark}
The notation $A(h,\varepsilon h)$ alone does not establish the existence of a cubic error
term. It might vanish if $C_2'(0)=0$, if symmetry forces cancellations, or if a different
mesh/area discretization fails the joint expansion assumption. Moreover, a fixed physical
displacement $\delta$ and an $h$-scaled displacement define different refinement families; only
comparisons made with identical geometric rules and numerical precision isolate this effect.
\end{remark}

\subsection{Scope of the numerical comparison}
The previously reported numerical comparison of fixed-map and scale-dependent
perturbations has not been independently reproduced from the original
implementation and is omitted from this release. The conditional
Proposition~\ref{prop:cubic} is an algebraic implication of its stated joint
expansion hypothesis, not a validation of a particular perturbation code.

\section{Generalized N-Series operators for known error exponents}
Suppose the approximation has an expansion
\[
A(h)=A+\sum_{p\in P}c_ph^p+\text{higher-order terms},
\]
where $P=\{p_1,\dots,p_m\}$ is a prescribed set of distinct positive exponents. Using $m+1$
refinement levels $h,h/2,\dots,h/2^m$, seek
\begin{equation}\label{eq:gen}
\Ahat(h)=\sum_{j=0}^{m}w_jA(h/2^j),\qquad
\sum_{j=0}^{m}w_j=1,\qquad
\sum_{j=0}^{m}w_j2^{-jp_k}=0\quad(k=1,\dots,m).
\end{equation}
The associated Vandermonde system has a unique solution for distinct exponents. If
$P=\{2,3\}$, the three-scale weights are
\begin{equation}\label{eq:P23}
\Ahat^{\{2,3\}}(h)=\frac{A(h)}{21}-\frac{4A(h/2)}{7}+\frac{32A(h/4)}{21}.
\end{equation}
For $P=\{2,3,4\}$ the four-scale formula is
\begin{equation}\label{eq:P234}
\Ahat^{\{2,3,4\}}(h)=-\frac{A(h)}{315}+\frac{4A(h/2)}{45}-\frac{32A(h/4)}{45}+\frac{512A(h/8)}{315}.
\end{equation}
For $P=\{2,3,4,5\}$ the five-scale formula is
\begin{equation}\label{eq:P2345}
\Ahat^{\{2,3,4,5\}}(h)=\frac{A(h)}{9765}-\frac{4A(h/2)}{651}+\frac{32A(h/4)}{279}
-\frac{512A(h/8)}{651}+\frac{16384A(h/16)}{9765}.
\end{equation}
For a term $c_ph^p$ the residual coefficient is multiplied by the exact moment
$M_p=\sum_{j=0}^{m}w_j2^{-jp}$. Direct rational evaluation gives
\begin{equation}\label{eq:moments}
M^{\{2,3\}}_4=\frac1{56},\qquad M^{\{2,3,4\}}_5=-\frac1{960},\qquad M^{\{2,3,4,5\}}_6=\frac1{31744}.
\end{equation}
Thus, when the indicated next coefficient is nonzero and a sufficiently long expansion holds,
the three formulas have leading residuals $c_4h^4/56$, $-c_5h^5/960$, and $c_6h^6/31744$,
respectively, and leading orders 4, 5, and 6. The sums of absolute weights are approximately
2.143, 2.429, and 2.585, respectively: these give modest worst-case amplification factors for
independent bounded input errors, although severe cancellation can still demand high working
precision. They require more refinement levels than the even-only operators to reach a
comparable order, and do not guarantee improved wall-clock efficiency. Numerical conditioning
and computational cost should be measured rather than presumed.

%------------------------------------------------------------
\section{Boundary irregularity, mesh family, and geometric validity}
The original polygon tests separated fixed irregular boundaries from changes in the mesh rule.
In a further boundary stress sweep, five-vertex cases at radii $15^\circ$ and $30^\circ$ with
jitter 0.2 and 0.5, and eight-vertex cases at jitter 0.2, retained near-$2,4,6,8$ behavior.
Eight-vertex jitter-0.5 cases did not converge to the same reference area, with raw errors
nearly constant and observed orders near zero. Such cases suggest a geometry/fan mismatch
(possibly orientation, nonconvexity, overlap, or a different limiting region), not a failure of
Richardson cancellation for an otherwise valid asymptotic expansion. For all test families, the
domain, orientation, target area, and approximation rule must be specified unambiguously.

%------------------------------------------------------------
\section{Limitations and validation priorities}
The retained original results are numerical except for the analytic cap parity argument,
the uniform latitude--longitude result in Proposition~\ref{prop:latlon-even}, and the
algebraic extrapolation identities. The original tests use convex or otherwise well-behaved
spherical domains with available reference areas. Arbitrary nonconvex, self-intersecting,
adaptive, and highly irregular meshes are outside the established scope. Proposition~\ref{prop:cubic} is
conditional: no joint expansion theorem for the original perturbed triangulation is proved
here. The numerical anomalies of the original perturbed-grid Table~6 are \emph{not} resolved by this revision;
the cubic mechanism is a candidate explanation only. Supplementary AI-assisted diagnostic fit tables requiring missing node and precision metadata are omitted from this release. To
elevate the proposed mechanism to a numerical result for this paper, one should (i) test
fixed-map and $h$-dependent-map meshes with the same node field, (ii) retain boundary
constraints and triangle orientation checks, (iii) sweep several $N$ and deformation
amplitudes, (iv) fit the $h^3$ coefficient directly, (v) verify precision stability, and
(vi) compare both error-versus-cost and observed-order-versus-cost across even-only and
mixed-power operators.

\subsection{Falsifiable cubic-coefficient test}
A sharper test of Proposition~\ref{prop:cubic} is the amplitude law, not merely a measured convergence order.
Choose a single smooth, boundary-vanishing perturbation vector field $V$ and a fixed chordal
triangulation convention. Here $h=N^{-1}$ in normalized cell coordinates (equivalently,
$h=H/N$ with $H=1$); $\delta=\varepsilon/N$ is the dimensionless displacement amplitude
multiplying $V$. At signed amplitude $\varepsilon$, displace interior parameter-grid nodes by
$\delta V$; retain the same physical domain and check triangle orientation. With
$e(h,\varepsilon)=A(h,\varepsilon h)-A$, fit, over several successively refined grids and at
sufficiently high precision,
\begin{equation}\label{eq:fit}
e(h,\varepsilon)=b_2(\varepsilon)h^2+b_3(\varepsilon)h^3+b_4(\varepsilon)h^4+\cdots .
\end{equation}
Under the joint expansion \eqref{eq:joint}, the precise small-amplitude prediction is
$b_3(\varepsilon)=\varepsilon C_2'(0)$, including $b_3(0)=0$; a nonzero slope requires
$C_2'(0)\neq0$. Test symmetric amplitudes $0,\pm\varepsilon_1,\pm\varepsilon_2$ with at least
two distinct nonzero magnitudes. Check that $[b_3(\varepsilon)-b_3(-\varepsilon)]/(2\varepsilon)$
stabilizes to one value, while the even-in-$\varepsilon$ part of $b_3$ tends to zero; compare the
inferred slope across refinement windows and increased precision.

By Proposition~\ref{prop:cubic}, the predicted cubic coefficients of the extrapolated errors are
\[
-\tfrac16\,\varepsilon C_2'(0)\ \ (\text{two-scale}),\qquad
-\tfrac1{90}\,\varepsilon C_2'(0)\ \ (\text{three-scale}),\qquad
-\tfrac1{810}\,\varepsilon C_2'(0)\ \ (\text{four-scale}).
\]
These coefficients are conditional predictions, not fitted measurements. A numerical test must distinguish failure of the joint expansion, vanishing derivatives, and pre-asymptotic effects; observed convergence orders alone cannot prove the hypothesis.

\section{Relation to the general N-Series framework}
Spherical area supplies another setting in which geometrically defined approximants can exhibit
extrapolatable error modes. This contribution extends the range of examples, not the
mathematical novelty of Richardson--Romberg weighting. In particular, a provable parity or
error-expansion theorem for a specified mesh/area rule is more informative than an isolated
high observed-order ratio.

%------------------------------------------------------------
\section{Conclusion}
For several fixed spherical refinement families, the reported experiments show the
$2,4,6,8$ error-order ladder, while for spherical caps and uniform latitude--longitude chordal cells (away from the poles)
the even-power structure follows analytically from the corresponding geometric symmetries. The perturbation results motivate a more
exact distinction: an irregular mesh may still be scale-consistent, while a scale-dependent
deformation can cause the expansion coefficients to vary with $h$ and generate odd powers. The
resulting third-order obstruction, when its coefficient is nonzero, is consistent with, and
would account for, the stalling of even-power cancellation in the perturbed experiments; its
residual factors are $-1/6$, $-1/90$ and $-1/810$ for the two-, three- and four-scale operators.
The numerical anomalies of Table~6 are not shown by this paper to be fully resolved. If the
error exponents are identified, generalized multiscale cancellation is available. The precise
analytical hypotheses and independent numerical verification remain central to any broader
spherical area law.

%------------------------------------------------------------
\appendix
\section{Reproducibility notes for the original experiments}
The original June 2026 manuscript identifies the following standalone Python scripts (names only; this TeX/PDF release does not include their executable contents):
\begin{itemize}
\item \texttt{spherical\_area\_nseries\_cap\_test.py}
\item \texttt{spherical\_triangle\_nseries\_area\_test.py}
\item \texttt{spherical\_latlon\_cell\_nseries\_area\_test.py}
\item \texttt{spherical\_irregular\_polygon\_nseries\_test.py}
\item \texttt{spherical\_irregular\_polygon\_batch\_nseries\_test.py}
\item \texttt{spherical\_irregular\_polygon\_stress\_sweep\_nseries\_test.py}
\item \texttt{spherical\_perturbed\_latlon\_nseries\_test.py}
\end{itemize}
For errors $e_N=|A_N-\Aex|$, the observed order is $p_N=\log(e_N/e_{2N})/\log2$; the
corresponding definition applies to the extrapolated errors. Some original experiment scripts were not included in the source/PDF review package, and their archived execution status is not established here. Numerical tables are attributed to the original manuscript; this statement is not a reproducibility certification.

\section{Revision history and provenance}
\textbf{Version 1 (June 3, 2026):} five spherical area tests, tables of observed orders, cap
parity argument, classical even-power scale operators, and limitations.

\textbf{Version 2 (October 24, 2026):} precise expansion-based regularity definition;
conditional $h^3$ perturbation proposition; AI-assisted exploratory diagnostics, distinctly
labeled; generalized mixed-power operators; stronger limits on what numerical observations
establish. The original numerical values are preserved rather than silently replaced.

\textbf{Version 2.1 (October 24, 2026):} corrected the regularity needed for the $O(h^5)$
remainder; clarified provenance as exploratory AI-assisted review; added first surviving exact
moments, error amplification factors, and a falsifiable signed-amplitude test of
cubic-coefficient linearity.

\textbf{Review supplement (October 24, 2026):} added specified deformation fields,
amplitude-fit sensitivity and provenance distinctions, plus an AI-assistance disclosure. The
original Section~9 perturbation remains unverified.

\textbf{Version 2.2 (October 25, 2026):} (1) stated the exact cubic residual factors
$-1/6$, $-1/90$, $-1/810$ in Proposition~1 (now Proposition~\ref{prop:cubic}) and Section~12.1, with a product-of-stages form of
the operators in Section~3; (2) Eq.~(6) numerator written as $s^2\sin h$ with a one-line
derivation; (3) removed the redundant condition $c\neq-1$ and noted that the constant term
$2\pi a_1$ equals the cap area; (4) the Section~6 triangle is described as birectangular, with
the reason the area equals $\beta$; (5) Section~9 and the Conclusion now describe the cubic
mechanism as consistent with, not established by, Table~6, and request a definition of what the
three seeds vary; (6) supplementary fit data with incomplete node, precision, and noise-floor provenance are omitted from this curated release;
(7) the shift operator $T$ is now used (product form of the operators in Section~3, unnumbered so that equation numbers match v2.1); (8) this entry. No numerical table values
were changed. Those historical fit-provenance issues are not part of the current results.

\textbf{Version 2.3 release candidate (October 29, 2026):} added
Proposition~\ref{prop:latlon-even}, an even-power asymptotic expansion for uniform
latitude--longitude chordal triangulations strictly between the poles. The
proof uses local diagonal parity, a uniform analytic area expansion, and
midpoint Euler--Maclaurin summation. Updated the abstract, scope and conclusion
without changing the original numerical tables or removing outstanding
provenance qualifications. The source/PDF pair is synchronized; supplementary numerical provenance limitations remain expressly disclosed.

%------------------------------------------------------------
\begin{thebibliography}{9}
\bibitem{coxeter} H. S. M. Coxeter, \emph{Introduction to Geometry}, Wiley, 1969.
\bibitem{todhunter} I. Todhunter, \emph{Spherical Trigonometry}, Macmillan, 1886.
\bibitem{richardson} L. F. Richardson, ``The approximate arithmetical solution by finite
differences of physical problems,'' \emph{Philosophical Transactions of the Royal Society A},
vol.~210, 1911, pp.~307--357.
\bibitem{williamson} D. L. Williamson, ``The Evolution of Dynamical Cores for Global
Atmospheric Models,'' \emph{Journal of the Meteorological Society of Japan}, vol.~85B, 2007,
pp.~241--269.
\end{thebibliography}

\end{document}

% Curated release: unsupported AI-assisted fit sections omitted, original experimental provenance remains qualified.